\subjecttitle{ARTIFICIAL INTELLIGENCE}{CT 710}{3}{1}{1.5}{IV}{I}

\cobjectives
  The main objectives of this course are: To provide basic knowledge of Artificial Intelligence, To familiarize students with different search techniques, To acquaint students with the fields related to AI and the applications of AI.
  \VS

\begin{enumerate}
  \item \textbf{Introduction \hfill (7 hours)}
    \begin{enumerate}
      \item Definition of Artificial Intelligence
      \item Importance of Artificial Intelligence
      \item AI and related fields
      \item Brief history of Artificial Intelligence
      \item Applications of Artificial Intelligence
      \item Definition and importance of Knowledge, and learning.
    \end{enumerate}
    \vs

  \item \textbf{Problem solving \hfill (7 hours)}
    \begin{enumerate}
      \item Defining problems as a state space search
      \item Problem formulation
      \item Problem types, Well-defined problems, Constraint satisfaction problem
      \item Game playing, Production systems.
    \end{enumerate}
    \vs

  \item \textbf{Search techniques \hfill (9 hours)}
    \begin{enumerate}
      \item Uninformed search techniques- depth first search, breadth first search, depth limit search, and search strategy comparison
      \item Informed search techniques- hill climbing, best first search, greedy search, A* search
      \item Adversarial search techniques- minimax procedure, alpha beta procedure
    \end{enumerate}
    \vs

  \item \textbf{Knowledge representation, inference and reasoning \hfill (4 hours)}
    \begin{enumerate}
      \item Formal logic-connectives, truth tables, syntax, semantics, tautology, validity, well-formed-formula
      \item Propositional logic, predicate logic, FOPL, interpretation, quantification, horn clauses
      \item Rules of inference, unification, resolution refutation system (RRS), answer extraction from RRS, rule based deduction system
      \item Statistical Reasoning-Probability and Bayes' theorem and causal networks, reasoning in belief network
    \end{enumerate}
    \vs

  \item \textbf{Structured knowledge representation \hfill (4 hours)}
    \begin{enumerate}
      \item Representations and Mappings
      \item Approaches to Knowledge Representation
      \item Issues in Knowledge Representation
      \item Semantic nets, frames
      \item Conceptual dependencies and scripts
    \end{enumerate}
    \vs

  \item \textbf{Machine learning \hfill (6 hours)}
    \begin{enumerate}
      \item Concepts of learning
      \item Learning by analogy, Inductive learning, Explanation based learning
      \item Neural networks
      \item Genetic algorithm
      \item Fuzzy learning
      \item Boltzmann Machines
    \end{enumerate}
    \vs

  \item \textbf{Applications of AI \hfill (14 hours)}
    \begin{enumerate}
      \item Neural networks
        \begin{enumerate}
          \item Network structure
          \item Adaline network
          \item Perceptron
          \item Multilayer Perceptron, Back Propagation
          \item Hopfield network
          \item Kohonen network
        \end{enumerate}

      \item Expert System
        \begin{enumerate}
          \item Architecture of an expert system
          \item Knowledge acquisition, induction
          \item Knowledge representation, Declarative knowledge, Procedural knowledge
          \item Development of expert systems
        \end{enumerate}

      \item Natural Language Processing and Machine Vision
        \begin{enumerate}
          \item Levels of analysis: Phonetic, Syntactic, Semantic, Pragmatic
          \item Introduction to Machine Vision
        \end{enumerate}
    \end{enumerate}
    \vs
\end{enumerate}
\Vs

\noindent \textbf{Practical:}
  \begin{enumerate}[label=\arabic*.]
    \item Laboratory exercises should be conducted in either LISP or PROLOG.
    \item Laboratory exercises must cover the fundamental search techniques, simple question answering, inference and reasoning.
  \end{enumerate}
  \Vs

\noindent \textbf{References:}
  \begin{enumerate}[label=\arabic*.]
    \item E. Rich and Knight, Artificial Intelligence, McGraw Hill, 2009.
    \item D. W. Patterson, Artificial Intelligence and Expert Systems, Prentice Hall, 2010.
    \item P. H. Winston, Artificial Intelligence, Addison Wesley, 2008.
    \item Stuart Russel and Peter Norvig, Artificial Intelligence A Modern Approach, Pearson, 2010
  \end{enumerate}

\newpage